\chapter{Literature Review}

\section{Introduction}

Modern game engines such as Unity, Unreal, or custom engines frequently relies on advanced mathematics to achieve various tasks like linear algebra for 3D rendering. However, the academic literature that
underlines these methods often presume a specific level of existing understanding in mathematics or physics. This creates a barrier-of-entry for hobbyists, indie developers, or anyone else who
lack knowledge in these areas.

The purpose of this literature review is to:
    \begin{enumerate}[label=(\roman*)]
        \item survey the current state of existing research of mathematical concepts often used in game development
        \item evaluate how effectively these papers communicate these concepts to a non-specialist audience
        \item identify any gaps that may justify the need for more practical (as in, opposed to practical)-orientated resources
    \end{enumerate}

\break

\section{Foundational Mathematics in Game Development}

\begin{table}[h]
    \centering
    \small
    \renewcommand{\arraystretch}{1.75}
    \caption{Mathematical Concepts Used in Game Development}
    \label{tab:game-maths-concepts}
        \begin{tabularx}{\textwidth}{@{} l X X @{}}
        \toprule
        \textbf{Concept} & \textbf{Use} & \textbf{Concepts Feature In} \\ \midrule
        Linear Algebra (Vectors, Matrices, Quaternions) & Object transformations, viewport manipulation & \bookcite{damQuaternionsInterpolationAnimation}{6} \bookcite{dunn3DMathPrimer2011}{} \\
        Differential Calculus (Derivatives, Gradients)  & Physics engines, procedural animation         & \bookcite{vossRealTimeInverseKinematics2025}{3}  \bookcite{koschierSurveySPHMethods2022}{} \\
        Probability and Statistics                      & AI decision-making, procedural generation     & \bookcite{snodgrassLearningGenerateVideo2017}{} \bookcite{dereszynskiLearningProbabilisticBehavior2011}{} \\
        Vector Calculus                                 & Particle simulations                          & \bookcite{koschierSurveySPHMethods2022}{}  \bookcite{wangFluidSimulationNavierStokes}{} \\ \bottomrule
    \end{tabularx}
\end{table}

Whilst being the most popular of these concepts, Linear Algebra's most `basic' operations such as matrix multiplications are often represented with \LaTeX{} notation that developers must manually
translate into code. Few papers, such as \bookcite{aversaUnityArtificialIntelligence2022}{} attempt to bridge this gap by providing code examples, however the paper still assumes a level of fluency
with the mathematical concepts that many indie developers may lack.

\section{Algorithms \& Accessibility}

\begin{table}[h]
    \centering
    \small
    \renewcommand{\arraystretch}{1.75}
    \caption{Algorithms and Accessibility in Papers}
    \label{tab:paper-algorithms-accessibility}
    \begin{tabularx}{\textwidth}{@{} l X X @{}}
        \toprule
        \textbf{Technique} & \textbf{Representation} & \textbf{Alternatives} \\ \midrule
        Procedural Generation & Noise functions are expressed through the use of summation notation (e.g., $\sum P(i)$) & Unity's Perlin Noise API. Documented resources such as tutorials (e.g., \bookcite{glazychevaUnity3DProceduralGenerated2020}{})  \\
        Physics Engine & Equations of motion expressed through the use of differential notation ($\partial^2 x / \partial t^2 = F/m$)  & Bullet, PhysX, Jolt. Most documentation omit the fundamental mathematics. \bookcite{JoltPhysicsJolt}{} \\
        AI Pathfinding & Expressed through Markov Chain ($\pi(\theta)$) and matrix notation & A* and Dijkstra's algorithms are very popular, and often omit the probability theory that make them operate \\ \bottomrule
    \end{tabularx}
\end{table}

Whilst academic papers present the mathematics in its most unadulterated form, communities have developed standalone libraries that abstract most of the complexity. However, these libraries are rarely
accompanied by the theoretical context beside it, which would help a developer understand \textit{how} an algorithm works, which may limit their ability to expand or optimize the code.

\section{Documentation \& Communication}

\begin{table}[h]
    \centering
    \small
    \renewcommand{\arraystretch}{1.75}
    \caption{Documentation and Communication Practices in Papers}
    \label{tab:paper-documentation-communication}
    \begin{tabularx}{\textwidth}{@{} l X X @{}}
        \toprule
        \textbf{Source} & \textbf{Strengths} & \textbf{Weaknesses} \\ \midrule
        Peer-Reviewed Journals & Very Comprehensive & Very dense in notation, little to no code examples \\
        Conferences & Primarily focuses on practical implementations, live demos & Lacks peer review, may over simplify the mathematics in sake of brevity \\
        Industry Blogs \& Tutorials & Code-focused, step-by-step explanations & Usually limited in theoretical explanation, and can be outdated \\
        Open-Source Projects & Very accessible code, community contribution / feedback & Documentation can be inconsistent or lacking \\ \bottomrule
    \end{tabularx}
\end{table}

These formats often excel at their strengths, but lack the structure to convey that which they lack. Due to this, it is quite hard to find a source that achieves to provide a structured
synthesis between the mathematical notation and code. This disconnect is what prevents the vast majority of the sources in the categories to be accessible.

\section{How Do We Determine If a Paper is Accessible?}

\subsection{Accessibility Metric System}

\subsubsection*{Definitions of Concept Approachability Dimensions}

We can determine approachability for any given paper through categorising each \textit{`dimension'} of approachability:

\begin{table}[h]
    \centering
    \small
    \renewcommand{\arraystretch}{1.75}
    \begin{threeparttable}
        \caption{Concept Approachability}
        \label{tab:conceptual-approachability}
        
        \begin{tabularx}{\textwidth}{@{} l X X @{}}
            \toprule
            \textbf{Dimension} & \textbf{Definition} & \textbf{Indicator Example} \\ \midrule
            Notation & To which degree are mathematical symbols explained or avoided & A `notation glossary' maybe present \\
            Clarity of Lexicon & To which degree is plain language use, and whether technical jargon is avoided & Number of technical terms in any given chapter / section of the paper \\
            Contextualisation & To which degree do papers provide examples, diagrams, or code snippets to help provide context to the topic discussed & At least one code snippet \\
            Pedagogy & To which degree do papers aim to teach the reader about a specific topic & Any attempt to help teach the reader what that concept or notation is \\ \bottomrule
        \end{tabularx}

        \begin{tablenotes}[para,flushleft]
            \small
            \textit{Note.}
            The difference between contextualisation and pedagogy lies in intent:
            contextualisation provides clarity to the reader on \textit{what} a concept is (e.g., through examples or diagrams), where as
            pedagogy involves and active attempt to \textit{teach} the reader how and why it works.
        \end{tablenotes}
    \end{threeparttable}
\end{table}

Each dimension will be scored on a 0-4 scale (0 = no accessibility, 4 = very accessible). The overall \textit{Approachability Score} (AS) will
be the mean value of the four sub-scores for that paper.

\begin{table}[H]
    \centering
    \small
    \caption{Literature Approachability}
    \label{tab:literature-approachability}
    \begin{tabularx}{\textwidth}{@{} >{\hsize=0.8\hsize\raggedright\arraybackslash}X >{\hsize=1.4\hsize\raggedright\arraybackslash}X @{}} 
        \toprule
        \textbf{Source} & \textbf{Contribution to Criteria} \\ \midrule
        \bookcite{aversaUnityArtificialIntelligence2022}{} & Features code comments, but expects prior knowledge of the basics of AI theory \\ \addlinespace
        \bookcite{damQuaternionsInterpolationAnimation}{} & Frequently uses advanced mathematical notation with very little explanatory text \\ \addlinespace
        \bookcite{koschierSurveySPHMethods2022}{} & Very dense advanced mathematical notation with few diagrams or explanatory text \\ \addlinespace
        \bookcite{vossRealTimeInverseKinematics2025}{} & Whilst still dense in mathematical notation, it provides some indicators of what some values and parts of equations mean \\ \bottomrule
    \end{tabularx}
\end{table}

These studies suggest that \textbf{notation} and \textbf{contextualisation} are the most frequently neglected dimensions. However, in contrast,
the \textit{lexical clarity} tend to be adequate in industry blogs but is rarely ever present in peer-review papers.

\subsection{Identified Gaps and What They Imply}

\begin{enumerate}
    \item \textbf{Notation Barrier: }
    \item \textbf{Lack of Code Examples: }
    \item \textbf{Lack of Pedagogy: }
\end{enumerate}

\pagebreak

\section{Questionnaire}

\subsection{Designing the Questionnaire}

The questionnaire will be informed by the gaps identified above. It will feature questions such as:

\begin{table}[H]
    \centering
    \small
    \caption{Questionnaire Example Questions}
    \label{tab:questionnaire-examples}
    \begin{tabularx}{\textwidth}{@{} >{\hsize=0.8\hsize\raggedright\arraybackslash}X >{\hsize=0.6\hsize\raggedright\arraybackslash}X @{}} 
        \toprule
        \textbf{Question} & \textbf{Rationale} \\ \midrule
        How comfortable are you translating mathematical equations (e.g., \(\partial^2 x / \partial t^2\)) into code? & Evaluates the `lack of code examples' gap \\ \addlinespace
        When reading academic papers, how often do you need to look up definitions for the symbols and concepts featured? & Evaluates the notation dimension described in table \ref{tab:conceptual-approachability} \\ \addlinespace
        Does the presence of pseudo\-code or code snippets increase your confidence in understanding a given algorithm or concept? & Validates the contextualisation dimension described in table \ref{tab:conceptual-approachability} \\ \addlinespace
        Do you find that visual diagrams (such as those that show coordinate spaces, or vector trajectories) are more informative than formal mathematical definitions when learning a new algorithm & Investigates the pedagogy dimension described in table \ref{tab:conceptual-approachability} \\ \addlinespace
        When using a physics or path-finding library, (e.g., Jolt, A*), how important is it to understand the underlying mathematical theory vs. just using the API? & Addresses the trade-offs mentioned in Section 3 regarding abstraction and theoretical contexts \\ \bottomrule
    \end{tabularx}
\end{table}

Each item will use a 5-point Likert scale (1 = Not at all comfortable, 5 = Very Comfortable). The questionnaire will also collect some demographic data (such as years of coding
experience, formal maths training) to determine any background variables that may influence the results.

\subsection{Scoring System}

For each paper examined in the literature review, each paper will be evaluated to determine how it approaches the four dimensions (Notation, Clarity of Lexicon, Contextualisation, Pedagogy)
using the 0-4 scale. The mean of these scores will constitute the \textbf{Approachability Score (AS)}:

\begin{equation}
    AS = \frac{1}{4} \sum_{d=1}^{4} S_d
\end{equation}

Where $S_d$ is the sub-score for dimension $d$.
The questionnaire data will then be cross-referenced with the AS. The papers with low scores should correspond to lower comprehension of the developers. This would validate the scoring system.
